% GENERATED FILE. Edit canonical lessons, then run npm run generate:books.
% canonical-source-hash: fnv1a64:1cfcf9eb4ebb3149
% canonical-lessons: GE-C181-fotografieren, GE-C181-fruehstuecken, GE-C181-giessen, GE-C181-gruessen, GE-C181-informieren

\chapter{To photograph, To have breakfast, To pour, To greet, To inform}
\label{ch:ge-to-photograph-to-have-breakfast-to-pour-to-greet-to-inform}
\hlchaptermodality{181}

\begin{chapteropening}
\textbf{By the end of this chapter:} \emph{I can say to photograph, to have breakfast, to pour, to greet and to inform.}

Complete the last lesson of chapter 181: I can say to photograph, to have breakfast, to pour, to greet and to inform.
\end{chapteropening}

\section[fotografieren]{fotografieren — to photograph}

\label{lesson:GE-C181-fotografieren}



\begin{quote}
Before the new one: say the German for to hold on to, then the German for to whisper.
\end{quote}

\subsection*{You'll want to know: fotografieren}
\textbf{fotografieren} — \textquotedblleft{}to photograph\textquotedblright{}.

\begin{grammarlens}[title={What you've built}]
One more word for this chapter.
\end{grammarlens}

\subsection*{Your turn}
\begin{itemize}
\raggedright
  \item \emph{Say it:} \emph{fotografieren}
  \item \emph{Say it:} \emph{fotografieren}, once more
  \item \emph{Say it:} say \emph{fotografieren}
  \item \emph{Recall:} say the German for to hold on to, then the German for to whisper, then say \emph{fotografieren} again
\end{itemize}

\begin{tcolorbox}[breakable,colback=teal!4,colframe=teal!35!black,title={Before you move on}]
What does fotografieren mean? (\textquotedblleft{}To photograph\textquotedblright{}.) Say it once more.
\end{tcolorbox}
\section[frühstücken]{frühstücken — to have breakfast}

\label{lesson:GE-C181-fruehstuecken}



\begin{quote}
Before the new one: say the German for to whisper, then the German for to photograph.
\end{quote}

\subsection*{You'll want to know: frühstücken}
\textbf{frühstücken} — \textquotedblleft{}to have breakfast\textquotedblright{}.

\begin{grammarlens}[title={What you've built}]
One more word for this chapter.
\end{grammarlens}

\subsection*{Your turn}
\begin{itemize}
\raggedright
  \item \emph{Say it:} \emph{frühstücken}
  \item \emph{Say it:} \emph{frühstücken}, once more
  \item \emph{Say it:} say \emph{frühstücken}
  \item \emph{Recall:} say the German for to whisper, then the German for to photograph, then say \emph{frühstücken} again
\end{itemize}

\begin{tcolorbox}[breakable,colback=teal!4,colframe=teal!35!black,title={Before you move on}]
What does frühstücken mean? (\textquotedblleft{}To have breakfast\textquotedblright{}.) Say it once more.
\end{tcolorbox}
\section[gießen]{gießen — to pour, to water}

\label{lesson:GE-C181-giessen}



\begin{quote}
Before the new one: say the German for to photograph, then the German for to have breakfast.
\end{quote}

\subsection*{You'll want to know: gießen}
\textbf{gießen} — \textquotedblleft{}to pour, to water\textquotedblright{}.

\begin{grammarlens}[title={What you've built}]
One more word for this chapter.
\end{grammarlens}

\subsection*{Your turn}
\begin{itemize}
\raggedright
  \item \emph{Say it:} \emph{gießen}
  \item \emph{Say it:} \emph{gießen}, once more
  \item \emph{Say it:} say \emph{gießen}
  \item \emph{Recall:} say the German for to photograph, then the German for to have breakfast, then say \emph{gießen} again
\end{itemize}

\begin{tcolorbox}[breakable,colback=teal!4,colframe=teal!35!black,title={Before you move on}]
What does gießen mean? (\textquotedblleft{}To pour, to water\textquotedblright{}.) Say it once more.
\end{tcolorbox}
\section[grüßen]{grüßen — to greet}

\label{lesson:GE-C181-gruessen}



\begin{quote}
Before the new one: say the German for to have breakfast, then the German for to pour.
\end{quote}

\subsection*{You'll want to know: grüßen}
\textbf{grüßen} — \textquotedblleft{}to greet\textquotedblright{}.

\begin{grammarlens}[title={What you've built}]
One more word for this chapter.
\end{grammarlens}

\subsection*{Your turn}
\begin{itemize}
\raggedright
  \item \emph{Say it:} \emph{grüßen}
  \item \emph{Say it:} \emph{grüßen}, once more
  \item \emph{Say it:} say \emph{grüßen}
  \item \emph{Recall:} say the German for to have breakfast, then the German for to pour, then say \emph{grüßen} again
\end{itemize}

\begin{tcolorbox}[breakable,colback=teal!4,colframe=teal!35!black,title={Before you move on}]
What does grüßen mean? (\textquotedblleft{}To greet\textquotedblright{}.) Say it once more.
\end{tcolorbox}
\section[informieren]{informieren — to inform}

\label{lesson:GE-C181-informieren}



\begin{quote}
Before the new one: say the German for to pour, then the German for to greet.
\end{quote}

\subsection*{You'll want to know: informieren}
\textbf{informieren} — \textquotedblleft{}to inform\textquotedblright{}.

\begin{grammarlens}[title={What you've built}]
One more word for this chapter.
\end{grammarlens}

\subsection*{Your turn}
\begin{itemize}
\raggedright
  \item \emph{Say it:} \emph{informieren}
  \item \emph{Say it:} \emph{informieren}, once more
  \item \emph{Say it:} say \emph{informieren}
  \item \emph{Recall:} say the German for to pour, then the German for to greet, then say \emph{informieren} again
\end{itemize}

\begin{tcolorbox}[breakable,colback=teal!4,colframe=teal!35!black,title={Before you move on}]
What does informieren mean? (\textquotedblleft{}To inform\textquotedblright{}.) Say it once more.
\end{tcolorbox}
